\section{Reduced finite-flat kernels and positive fibre maxima}
\textit{Supports: G331, n1.}

Let $R\to B$ be finite locally free of rank $d>0$, with both rings reduced. Over the local base choose a basis. Multiplication by $h=\sum T_jg_j$, with $g_j$ finite $R$-module generators of a $B$-ideal $J$, has characteristic polynomial
\[
\det(t-m_h)=t^d+\sum_{i=1}^d(-1)^ic_i(T)t^{d-i}.
\]
Let $C_i(J)$ be the coefficient ideal of $c_i(T)$ and $N(J)=C_d(J)$. Changing basis conjugates the matrix. Changing generators substitutes linear forms in $T$ in both directions, proving ideal independence. The determinant on the top exterior power proves arbitrary base-change compatibility; replacing $h$ by a product proves norm multiplicativity. This is the full immediate Stacks 0BD2 CURRENT 31.18.6 construction at a finite locally free morphism; its exterior-power proof was checked. For $J\supset\mathfrak a^NB$, add base elements $x\in\mathfrak a^N$ redundantly. The pure coefficient $x^d$ is in $N(J)$, so $N(J)$ contains a boundary power and is unit-or-primary.

For ANY base support $P$, first enter the valued field $K=\operatorname{Frac}(R/P)$. The fibre algebra $B\otimes_RK$ can be NONREDUCED even though $B$ is reduced. In a finite normal field extension splitting its residue fields, it becomes Artinian local factors of lengths $l_j$, residue equal to that field, with $\sum l_j=d$. Multiplication on a factor is its residue scalar plus a nilpotent map; thus
\[
\det(t-m_h)=\prod_j(t-h_j)^{l_j}.
\]
Generic inseparability and all special-fibre nilpotent lengths are included. Extend the base valuation to the integral closure in the splitting field and evaluate the residue homomorphisms $B\to K_j$ there. These give centred lifts $w_j$, and
\[
v(N(J))=\sum_j l_jw_j(J).                              \tag{3}
\]
Indeed the minimum coefficient value of a polynomial is multiplicative: quotient its support, divide by a coefficient of least finite value, and reduce to nonzero residue polynomials whose product is nonzero. Apply this Gauss argument to the determinant factorization. Boundary powers ensure finite ideal values here.

The measures $\kappa_v=\sum l_j\delta_{w_j}$ have mass $d$. Their integrals on every bounded ideal gauge are fixed by (3). The ideal-difference lattice is dense in ALL real continuous source functions: it contains constants, separates points via $(b)+\mathfrak a^N B$, and has the identity
\[
\min\{v(I)-v(J),v(K)-v(H)\}=v(IH+KJ)-v(JH).
\]
Center separation is by semilocal boundary ideals. Compact lattice approximation proves density (G331), so equal masses imply choice independence on all tests and hence as measures. Conversely any lift has a support over $P$; its residue field embeds into a splitting field, and extending that particular lift makes it appear in the Dirac sum with positive length. Choice independence therefore proves FULL fibre support and fibre cardinality at most $d$. First (3), then density with fixed mass, proves continuity on every continuous test. The openness criterion in n1 applies.

Define $\tau H(v)=\int H\,d\kappa_v$. It sends $-w(J)/q$ to $-v(N(J))/q$, preserves $C$, and sends an actual decreasing $C$-representation to a decreasing $C$-representation of the trace. A finite positive sum commutes with decreasing nonpositive extended nets: for finite limits approximate each of finitely many terms at one common index; if one limit is bottom its positive weight forces bottom. Thus $\tau$ preserves $D$, including partial bottom. No integration theorem for an arbitrary measurable net is used.

There is also a direct positive maximum. Put $m=\min_{\rho w=v}w(J)$. Then
\[
m=\min_{1\le i\le d}v(C_i(J))/i.                      \tag{4}
\]
All roots of multiplication by a generator have value at least $m$, so each degree-$i$ coefficient has value at least $im$, proving one inequality. Choose a generator attaining $m$ at a fibre point. Evaluating its characteristic equation there, if every nonleading coefficient had value $>im$, the leading root power would be uniquely least-valued, contradicting the equation. Hence some coefficient value is at most $im$. A specialized coefficient is an $R$-linear combination of the family coefficients, proving the reverse inequality in (4). This proof tolerates cancellation and inseparable roots.

The ideals $C_i$ need not be primary. Let $L=\sup_X v(C_d)/d<\infty$, choose an INTEGER $N>dL$, and replace every $C_i$ by the genuinely unit-or-primary $C_i+\mathfrak a^N$. Values are truncated at $N$; because $m\le L$ and $N/i>L$, the minimum in (4) is unchanged. Therefore
\[
\rho_{\max}(-w(J)/q)=\max_i\{-v(C_i+\mathfrak a^N)/(qi)\}\in C_R. \tag{5}
\]
The maximum is 1-Lipschitz in uniform norm, so it preserves $C$. On a finite fibre, maximum commutes with a decreasing net: if each limiting value is below a finite threshold, directedness gives one index with every value below it. Thus it preserves ACTUAL $D$ limits as well.

For a normal-domain target without vertical components, the same maximum theorem follows independently from G331's kernels. On each domain quotient take the integral closure in a finite normal splitting field. The maximum of the finitely many embedding pullbacks of a $C$ or $D$ source function belongs to that class upstairs. The full-support kernel shows that its value is exactly the full fibre maximum and is independent of the splitting-field lift. G331's positive FULL-fibre descent puts that maximum downstairs in $C,D$. Take the maximum across the finitely many dominating source components. Every point is covered, so this is the full maximum. Finite-fibre decreasing passage is exactly the preceding argument. This proves both openness AND geometric maximum preservation in the normal/no-vertical and reduced finite-flat cases; topology alone would not prove positivity.
